\section{开放 Auto Research：去中心化的智能改进}

\subsection{不可信工作者池}

Karpathy 正在探索如何让互联网上的\textbf{不可信工作者池}（untrusted pool of workers）参与 auto research。核心洞察：

\begin{importantbox}{搜索昂贵，验证廉价}
LLM 训练改进具有类似 blockchain 的性质：
\begin{itemize}
    \item 找到一个好的代码 commit 可能需要尝试 10,000 个想法（搜索很昂贵）
    \item 但验证一个 commit 是否有效只需训练一次模型（验证很便宜）
    \item 类似 SETI@Home、Folding@Home 的分布式计算范式
\end{itemize}
"Auto Research at Home"——互联网上的 agent 群体可以协作改进 LLM，甚至有可能超越 Frontier Labs。
\end{importantbox}

\subsection{FLOP 作为新货币}

Karpathy 抛出一个有趣的思考：未来人们关心的可能不是美元，而是 FLOP。"How many FLOPs do you control instead of what wealth do you control?" 虽然他自己也认为这不完全成立，但这个思想实验值得玩味。

\subsection{本章小结}

开放 auto research 将分布式计算与 AI 研究结合，通过"搜索昂贵、验证廉价"的性质，使不可信工作者池也能安全参与 LLM 改进。


